\documentclass[11pt,a4paper]{article}

% to be compiled with lualatex

%% begin comment out for tth
\usepackage{polyglossia}
\setdefaultlanguage{english}
%% end comment out for tth

\usepackage{parskip}
\usepackage[a4paper,hmargin=2.4cm,bottom=3.2cm]{geometry}
\usepackage{makeidx}
\usepackage{url}
\usepackage{graphicx}
\usepackage[x11names,svgnames]{xcolor}
\usepackage{enumitem}
%% \usepackage[para,ragged]{footmisc}

\usepackage{microtype}
% it may be necessary to comment out `microtype' with polyglossia loaded

% \usepackage{luaotfload}
% \usepackage{fontspec}
% \usepackage[default]{fontsetup}
\usepackage[olddefault]{fontsetup}
% \usepackage[times]{fontsetup}

\makeindex

%%  * for a run of tth (generate .html): 
%% - comment out "noitemsep" temporarily !!
%% - comment out \usepackage{polyglossia} and \setdefaultlanguage{english}
%% 


% Commands for indication changes in updated versions:
% changed, inserted
% \newcommand{\changedbegin}[1]{\textcolor{blue}{\fbox{\texttt{#1$>$}}}}
% \newcommand{\changedend}[1]{\textcolor{blue}{\fbox{\texttt{$<$#1}}}}
% \newcommand{\changedbeginA}[1]{\textcolor{DarkGreen}{\fbox{\texttt{#1$>$}}}}
% \newcommand{\changedendA}[1]{\textcolor{DarkGreen}{\fbox{\texttt{$<$#1}}}}


\usepackage[%
  bookmarksnumbered=true,
  colorlinks=true%
]{hyperref}

\usepackage[nottoc]{tocbibind}

\title{\textsf{elpres} --- electronic presentations \\ with (pdf/Lua)\LaTeX}
\author{Volker Kiefel\thanks{volker dot kiefel at freenet dot de, 
       \url{http://vkiefel.de/elpres.html}}}
\date{v1.1 \\[1.8ex]
  May 29, 2024%
  % --- \textcolor{DarkGreen}{update of this manual:
  %   September 27, 2020}
}




\begin{document}
\maketitle
\thispagestyle{empty}

\tableofcontents


% \section*{A note on this updated version of the \texttt{elpres} manual}

% \begin{sloppypar}
% In this updated version of the manual
% paragraphs or sentences with changes until August 19, 2020
% (compared to the ``official''
% manual for v0.6) are marked
% \changedbegin{changed} changed content \changedend{changed}
% and new text is marked \changedbegin{inserted} new content
% \changedend{inserted}. Additional changes until September 27, 2020:
% \changedbeginA{inserted} new text \changedendA{inserted},
% \changedbeginA{changed} modified text \changedendA{changed}.
% \end{sloppypar}

\section{Introduction}

\label{sec:Introduction}

The \texttt{elpres} class is intended to be used for
presentations on a computer screen, a
beamer or a projector. It is derived from \LaTeX's \texttt{article} class and
may be used with \LaTeX{}, pdf\LaTeX{} (and Lua\LaTeX ).
The default ``virtual paper size'' of slides produced by this class:
width=128mm, height=96mm which corresponds to a 4:3 (width:height) aspect
ratio.
\index{aspect ratio|see{screen aspect ratio}}%
\index{screen aspect ratio}%
Other aspect ratios for widescreen monitors may be selected by class
options.
The \texttt{elpres} class requires that the \texttt{ifthen},
\texttt{fancyhdr},
\texttt{hyperref},
\texttt{graphicx},
\texttt{xcolor}
and \texttt{geometry} packages are available on the \TeX{} system: these
packages are loaded automatically by the \texttt{elpres} class.
Enhancements to \texttt{elpres} are easily made available by other packages,
these include overlay support for incremental slides (package
\texttt{overlays})  and slides
with a background from a bitmap (\texttt{wallpaper}, \texttt{eso-pic}
packages). Predefined color/layout schemes for \texttt{elpres} presentations
can be acitivated with
\verb+\usepackage+ (details can be found in
in section~\ref{sec:pres-scheme})

This manual is intended to support the user with ``recipes''.  Use of
\texttt{elpres} with its default settings should be simple, additional aspects
including overlay functions, use of colors, graphics files, ``handout
documents'' are described in section \ref{secEnhancements}. Many code snippets
have been included in this manual.
They can be used in users' presentation files.

Some extensions described in this manual work only with \texttt{pdf}-files
which should preferrably be compiled with pdf\LaTeX\ or Lua\LaTeX.
\index{Lualatex@Lua\LaTeX}
\index{pdflatex@pdf\LaTeX}

\section{Installation}

\label{sec:installation}


\begin{sloppypar}

  \index{installation|(}
  If the \texttt{elpres} package has already been
  installed with the \TeX-system nothing needs to be done.  If an updated
  version shall be installed (or if the preinstalled \texttt{elpres} version
  does not work properly\footnote{for details see the description on checking the
    \texttt{elpres}-installation in section
    \ref{sec:CheckInstallation}}),
  \texttt{elpres.cls}, all \texttt{*.sty},
  \texttt{*.png} and \texttt{*.eps} files should be copied into a directory,
  where the \TeX-system can find it: if an old \texttt{elpres} version of the
  existing \TeX\ system shall be replaced by the current version, please copy
  these into the appropriate position in the ``local'' \texttt{texmf}
  directory tree (e.\,g. \verb+~/texmf/tex/latex/elpres+) and this manual
  (\texttt{elpres-manual.pdf}) to \verb+~/texmf/doc/latex/elpres+, where
  \texttt{~/texmf} may be \texttt{~/texmf-local} or an
  analogous directory.

Then, the files database should be updated by entering (in case of a \TeX~Live
installation) the command:

\end{sloppypar}

\begin{verbatim}
   mktexlsr ~/texmf-local
\end{verbatim}

or you may enter \verb+texhash ~/texmf-local+\footnote{The location
  for the ``local'', manually updated class and package files may be different,
  depending on the \TeX\ system and the OS.}.
Other \TeX-distributions e.\,g.\ MiK\TeX\ have
their own package updating mechanisms. The elpres manual file should be
accessible with the command ``\mbox{\texttt{texdoc elpres}}''. If this
still calls the old version of the manual, the command ``\texttt{texdoc -l elpres}''
will prompt you to select either the old or the new version of the manual.
\index{installation|)}

\section{Usage}
\label{secUsage}

\subsection{Class options, commands and environments}
\label{sec:EpOptCmdEnv}

The class is used with 
\begin{verbatim}
   \documentclass[options]{elpres}
\end{verbatim}

Slides are written with the \texttt{psli} and \texttt{rsli} environments.
This section lists all \texttt{elpres}-specific options, commands and
environments.

\subsubsection*{\texttt{Elpres}-specific options}

\index{elpres@\texttt{elpres}!options}
For font selection (\LaTeX, pdf\LaTeX):
 ``\texttt{tmrfont}'' (Times Roman), ``\texttt{helvetfont}''
(Helvetica), ``\texttt{cmfont}'' (Computer Modern) and ``\texttt{sansfont}''  
(Sans Serif: default) may be used. The option ``\texttt{nofonts}'' 
(no font selection) 
is intended for use of \texttt{elpres} with Lua\LaTeX\
\index{Lualatex@Lua\LaTeX}
with its own font selection mechanisms, an example is described in
section \ref{sec:MiniExample}.

\index{slide dimensions}
Options for different screen aspect ratios are 
``\texttt{4x3}'' (default),\footnote{slide size \texttt{128\,x\,92\,mm}} 
``\texttt{16x9}'',\footnote{slide size \texttt{176\,x\,99\,mm}} 
``\texttt{16x10}''.\footnote{slide size \texttt{160\,x\,100\,mm}} 
\index{bulletsymb class option@\texttt{bulletsymb} class option}
The option ``\texttt{bulletsymb}'' selects \verb+$\bullet$+ symbols (closed
circles)
instead of the default \texttt{itemize}-environ\-ment
symbols in all four levels,
``\texttt{ballsymb}'' is effective together with the presentation
schemes (section \ref{sec:pres-scheme})
and inserts small images of colored \emph{balls} as alternative 
symbols for the \texttt{itemize} environment. 
%
The option ``\texttt{navisymb}''
adds a small field with symbols for navigation to the right
bottom area of slides (section \ref{sec:NavigationSymbols}). 
Options of the \texttt{article} class are also available for
\texttt{elpres} presentations,
e.\,g.  \texttt{10pt}, \texttt{11pt}, \texttt{12pt} for
selection of font size. 

\subsubsection*{\texttt{Elpres}-specific commands}

{ \sloppy
\index{elpres@\texttt{elpres}!commands}
The command
``\verb+\distance+''
supports the user to vertically adjust lines and/or pragaraphs on a 
slide. (details are described in section 
section~\ref{sec:EssentialElements}).
The ``\verb+\auvimm+''\footnote{for \textbf{au}dio \textbf{vi}deo 
  \textbf{m}edia}
command inserts a link to an external audio or video file 
(section~\ref{sec:MultimediaContent}).
Two commands: ``\verb+\fromlinktext+'' and ``\verb+\totargettext+''%
\index{fromlinktext@\verb+\fromlinktext+}
\index{totargettext@\verb+\totargettext+}
can be used for internal links within a presentation 
(section~\ref{secHypertext}). 
With ``\verb+\slidetitlecolor+'', the text color of the title of slides 
(\texttt{psli}) can be changed (section \ref{sec:footerlayout}).
With the command ``\verb+\pagenrconst+'', the page number of
the current slide can be set back to the page number of the previous slide
(section \ref{secOverlay}, last paragraph). 
\par}

\subsubsection*{\texttt{Elpres}-specific environments}

\index{elpres@\texttt{elpres}!environments}
The environments ``\texttt{psli}'' and ``\texttt{rsli}'' define 
\emph{slides} in \texttt{elpres} (section \ref{sec:EssentialElements}); 
``\texttt{citemize}'',``\texttt{cenumerate}''
 and ``\texttt{cdescription}'' provide vertically compressed lists
(section \ref{sec:CompressedLists}).
\index{slides in elpres@slides in \texttt{elpres}}

\subsection{A minimalistic example presentation}

\label{sec:MiniExample}

The following code (\texttt{minimal.tex})
may be used for first ``experiments'' by adding
features described in later sections.

\index{minimalistic example, presentation}
\index{simple example, presentation|see{minimalistic example, presentation}}

\begin{verbatim}
   \documentclass[12pt,pdftex,4x3]{elpres}
   \usepackage[utf8]{inputenc}
   \usepackage[document]{ragged2e}

   % \usepackage{elpreswhitebluescheme}
   % \usepackage{elpresbluelightgrayscheme}
   % \usepackage{elpresgrayscheme}
   % \usepackage{elpreswhiteredscheme}
   \usepackage{elpreswhitetealscheme}

   \begin{document}
   \begin{titlepage}
   \centering
   \distance{1}
   {
   \Huge \bfseries \textcolor{eptitlecolor}{Title of the presentation} \par
   }
   \vspace{1.3ex} \large
   Author\\[2ex]Institution
   \distance{2}
   \end{titlepage}

   \begin{psli}[Title of Page]
   The first page: \texttt{psli} environment

   \begin{itemize}
   \item first line in an itemized list
   \item second line in an itemized list
   \end{itemize}
   \end{psli}

   \begin{rsli}
   The second page: \texttt{rsli} environment

   \begin{enumerate}
   \item no title
   \item not centered vertically

   \end{enumerate}

   \end{rsli}
   \end{document}
\end{verbatim}




The preamble of the same presentation for Lua\LaTeX\ would read:

\begin{verbatim}
   % LuaLaTeX: Please use utf-8 encoding!
   \documentclass[12pt,nofonts,4x3]{elpres}
   \usepackage[document]{ragged2e}
   \usepackage{fontspec}
   
   \setmonofont{TeX Gyre Cursor}              %% based on 'Courier'
   \setsansfont[Scale=0.92]{Tex Gyre Heros}   %% similar to 'Helvetica'
   \setmainfont{TeX Gyre Termes}              %% 'Roman' style serif font
   \renewcommand{\familydefault}{\sfdefault}
   
   % \usepackage{elpreswhitebluescheme}
   \usepackage{elpreswhitetealscheme}
   % \usepackage{elpresbluelightgrayscheme}
   % \usepackage{elpresgrayscheme}
   % \usepackage{elpreswhiteredscheme}
   
   \begin{document}
   % ...
   \end{document}
\end{verbatim}

\index{Lualatex@Lua\LaTeX}
The use of Lua\LaTeX\ with \texttt{elpres} is heavily recommended
due to the superior font selection mechanisms.
%% You may copy the code examples from this manual more easily from the
%% \texttt{.html} version of this manual available at
%% \url{http://vkiefel.de/elpres/elpres-manual.html}.

\subsection{Testing of the \texttt{elpres} installation}

\label{sec:CheckInstallation}

\begin{sloppypar}
With this example you may check, if \texttt{elpres} has been installed
correctly in your \TeX -system: copy the listing of \texttt{minimal.tex} into
a file, add the class option \texttt{ballsymb}:

\begin{verbatim}
   \documentclass[12pt,pdftex,4x3,ballsymp]{elpres}
   ...
\end{verbatim}



and process it with \texttt{pdflatex}. The line
\verb+\usepackage{elpreswhitetealscheme}+ should not be commented out,
i.\,e. the package \texttt{elpreswhitetealscheme} should be active.
If processing of \texttt{minimal.tex} results in a message like
\end{sloppypar}

\begin{verbatim}
  ... `ep-ball-04' not found ...
\end{verbatim}


the \TeX -system does not find the file \texttt{ep-ball-04.png} in the
appropriate place. In such a situation it is recommended 

\begin{enumerate}
  \item either to install the current version of \texttt{elpres} in the
    correct subdirectories of the ``local'' \texttt{texmf}-directory, details
    are described in section \ref{sec:installation}
  \item or to copy the small elpres graphics files \texttt{ep-ball-nn.eps} 
    \texttt{ep-ball-nn.png} into a directory of the \texttt{tex}-branch of the
    local  \texttt{texmf}-directoty tree (followed by an update
    see of \TeX s files database, for details see section 
    \ref{sec:installation}). The files may be obtained
    from the \texttt{elpres} project page 
    \url{http://vkiefel.de/elpres.html}%
    \footnote{\url{http://vkiefel.de/elpres/elpres-ball-image-files.zip}}
   \item or to copy the required image file into the current project
   directory.
\end{enumerate}

The problem sometimes arises as soon as \TeX-distributions do not install
\texttt{elpres} correctly. For details see \texttt{README.md} in the archive
\texttt{elpres.zip}, which may be obtained from \texttt{ctan.org}.

\subsection{Essential elements of an \texttt{elpres} presentation}
\label{sec:EssentialElements}

\index{title page}
The \textbf{title page} slide can be created with
the \texttt{titlepage}  environment (or \texttt{rsli}, see below), \LaTeX's 
\verb+\maketitle+ command is not available.
\index{psli environment@\texttt{psli} environment}
\textbf{Slides} may be created with the
\texttt{psli}-environment\footnote{\texttt{psli}: \textbf{p}lain
  \textbf{sli}de}, you may add
the title of the slide with the optional parameter. 
The contents of the slide are centered vertically:

\begin{verbatim}
   \begin{psli}[Slide title]
      Contents of slide
   \end{psli}
\end{verbatim}

\index{rsli environment@\texttt{rsli} environment}
Another environment generating a \textbf{slide} is
\texttt{rsli}\footnote{\texttt{rsli}: \textbf{r}aw \textbf{sli}de}:
slides are written 
without title, contents are not centered vertically:

\begin{verbatim}
   \begin{rsli}
      Contents of slide
   \end{rsli}
\end{verbatim}

\index{distance@\verb+\distance+}
The \verb+\distance{number}+ command allows to introduce vertical space into
slides constructed with the \texttt{rsli} and \texttt{titlepage} environments.
You may use of \verb+\distance{number}+ commands with numbers indicating
the relative height of empty spaces between lines
and/or paragraphs, see the title page in the example above (section
\ref{sec:MiniExample}).

The use of footnotes on slides is often problematic, if they cannot be
avoided, the \texttt{footmisc} package is recommended: the \texttt{perpage}
option resets numbering for each new slide. For a presentation, the
\texttt{symbol} option allows to use symbols instead of numbers\footnote{
 The use of the \texttt{footmisc} package may produce the \LaTeX\ error 
   message ``\texttt{Counter too large}''.
   Therefore it is preferred to use the \texttt{symbol*} and \texttt{perpage}
   options, details are described in the manual of the \texttt{footmisc}
   package.}.  
After inserting a new footnote, numbers or symbols are correctly
inserted only after a second run of \LaTeX.

\subsection{Alternatives to \TeX's justified paragraphs}

By default, \LaTeX\ produces justified paragraphs with lines of equal length,
this will often not be appropriate for the usually very short lines of text
in presentations. The \LaTeX\ \verb+\raggedright+ command
\index{raggedright@\verb+\raggedright+}
has its own
deficiencies: by inhibiting hypenation in texts with rather short lines,
the right margin will often look too
ragged. A solution is to use the \verb+\RaggedRight+ command of the
\texttt{ragged2e} package. 

If authors of presentations wish to suppress both justification and
hyphenation within the complete presentation document, the following
package parameters:

\begin{verbatim}
   \usepackage[document,originalparameters]{ragged2e}
\end{verbatim}

can be selected. As hypenated words are not favorable for quick reading of
slides by the audience this is the preferred use of \texttt{ragged2e} for
the author of this package.

\subsection{Vertically compressed lists}
\label{sec:CompressedLists}

As the spaces between lines may be too great with the \texttt{itemize}
environment,
the \texttt{elpres} package provides a ``vertically compressed''
\texttt{citemize}-environment:
\index{citemize environment@\texttt{citemize} environment}

\begin{verbatim}
   \begin{citemize}
     \item  one
     \item  two
   \end{citemize}
\end{verbatim}

Similarly, a \texttt{cenumerate} and a \texttt{cdescription} environment may
be used.
\index{cenumerate environment@\texttt{cenumerate} environment}
\index{cdescription environment@\texttt{cdescription} environment}

Another solution for the customization of \texttt{itemize}
environments is given by the \texttt{enumitem} package. Therefore

\begin{verbatim}
   \usepackage{enumitem}
\end{verbatim}

should be added to the preamble, and  a comma-sperated list of parameters
parameters can be added in the format:

\begin{verbatim}
   \begin{itemize}[parameter-list] 
      ... 
   \end{itemize} 
\end{verbatim}

The ``vertically compressed'' list can then be obtained with

\begin{verbatim}
   \begin{itemize}[nosep]
     \item  one
     \item  two
   \end{itemize}
\end{verbatim}

Similarly, the \texttt{enumitem} package is also able to modify the
\texttt{enumerate} and description \texttt{environments}.



\subsection{Slide layout: changes in slide title, footer, page numbers}

\label{sec:footerlayout}

By default, the \textbf{text color of the titles} of
\texttt{psli}-slides is black,
\index{slidetitlecolor@\verb+\slidetitlecolor+}
\index{title, text color}
it may be changed by redefining the ``value'' of the \verb+\slidetitlecolor{}+
command like
\index{slidetitlecolor@\verb+\slidetitlecolor+}

\begin{verbatim}
   \slidetitlecolor{blue}
\end{verbatim}

in the preamble. Of course, you may also use the named colors of the
\texttt{xcolor} package (see section \ref{sec:use-colors}).

\index{page number}
By default the \textbf{page number} appears at the right bottom of the slide
(in the ``footer'' in the terminology of the \texttt{fancyhdr} package)
this position can be addressed by the \verb+\rfoot{}+ command: the
\verb+\rfoot{}+, \verb+\cfoot{}+ and \verb+\lfoot{}+ commands can be used
by the author of a presentaton.

\begin{sloppypar}
However, access to the header fields defined by \texttt{fancyhdr}
is blocked: if you enter
  \verb+\rhead{}+, \verb+\chead{}+
  or \verb+\lhead{}+, \texttt{elpres} will issue
  an error message. Use of the header fields will generate unwanted effects
  on page layout due to \texttt{fancyhdr}. 
If you wish to change the position of the page number in the footer,
you can overwrite the default page number with an ``empty'' \verb+\rfoot{}+
command and put it to the center or the left margin of the footer
(\verb+\cfoot{\footnotesize \thepage}+ or 
\verb+\lfoot{\footnotesize \thepage}+). These ``footer fields''
\index{footer fields on slides}
are also suitable to enter a \textbf{logo}
\index{logo graphics file}
visible on all pages (in form of a graphics file using
\verb+\includegraphics[]{}+) or text with the name of the speaker's
institution.
\end{sloppypar}

The user may also use predefined \textbf{presentation schemes} with defined
colors and layout elements (e.\,g. symbols used in list envoronments), details
are described in section \ref{sec:pres-scheme}.



\section{Enhancements to \texttt{elpres}}

\label{secEnhancements}

Enhancements described here make use of other \LaTeX\ packages which should
be available in \TeX\ distributions like \TeX~Live or MiK\TeX.

\subsection{Use named colors of the \texttt{xcolor} package}
\label{sec:use-colors}

The \texttt{elpres} class automatically loads the \texttt{xcolor} package.
Therefore, color names defined by the \texttt{xcolor} may be used. They are
however only accessible in groups (\texttt{dvipsnames}, \texttt{svgnames},
\texttt{x11names}, details can be found in the \texttt{xcolor} manual). As
an example, the color \texttt{Indigo} is available in the \texttt{svgnames}
goup.  If you wish to use it you will have to enter \texttt{svgnames}
as \texttt{elpres} class option:

\begin{verbatim}
   \documentclass[11pt,16x9,svgnames]{elpres}
\end{verbatim}

This option is then automatically ``handed over'' to the \texttt{xcolor}
package loaded by \texttt{elpres}. The complete lists of named colors
are found in the \texttt{xcolor} manual.


\subsection{Include graphics files}
\label{secUsageGraph}

\index{graphics files}
\index{includegraphics@\texttt{\textbackslash includegraphics}}
Graphics files/pictures can be included with the
\texttt{includegraphics}-command of the \texttt{graphicx}-package. Please be
aware that the dimensions of the pages are \mbox{128mm x 96mm} and
therefore included graphics are scaled appropriately. A typical example
of a 
page with a picture could be (\texttt{pict.png} is the name of the
graphics file): 


\begin{verbatim}
   \usepackage[pdftex]{graphicx} % (in preamble) 
   ...
   \begin{rsli}
     \centering
     \distance{1}
     \includegraphics[width=0.9\textwidth,%
                      height=0.9\textheight,%
                      keepaspectratio=true]{pict.png}
     \distance{1}
   \end{rsli}
\end{verbatim}


The \verb+\includegraphics[]{}+ comannd requires to select the correct device
driver related option (e.\,g. \texttt{pdftex} or \texttt{dvips})
(documentclass).


% \label{secInclGraphicsFiles}


\subsection{Arrange text and pictures in two (or more) columns}

\label{secTwoColumns}

\index{columns}
\index{minipage environment@\texttt{minipage} environment}
Text and graphics may be arranged in two or more columns with
\texttt{minipage} environments:

\begin{verbatim}
   \begin{minipage}[b][0.8\textheight][t]{0.5\textwidth}
     \colorbox{white}{%
       \includegraphics[width=0.9\textwidth]{graphics-file.png}}
   \end{minipage}
   \begin{minipage}[b][0.8\textheight][t]{0.48\textwidth}
   \footnotesize
   \begin{citemize}
      \item ...
      \item ...
      ...
   \end{citemize}
   \end{minipage}
\end{verbatim}

Details on the minipage environment may be found in the \LaTeX{}
documentation\footnote{A good reference is \url{https://latexref.xyz/}}.


\subsection{Add incremental slides (overlays) to a presentation}

\label{secOverlay}

\index{incremental slides|see{overlays}}
\index{overlays|(}
If the contents of slides are to be made visible step
by step this can be achieved by a series of output PDF or (PS) pages (carrying
the same page number) usually called \emph{overlays}. It may also be of
interest to change a highlighting color in a series of overlays. This is most
easily done by using the excellent \texttt{overlays} package written by
Andreas Nolda. 

To generate a series of four overlays sequentially showing four lines of a
list:

\begin{enumerate}
  \item load the \texttt{overlays} package in the preamble
  \item put a \texttt{psli} or \texttt{rsli} slide environment into an
  \texttt{overlays} (or \texttt{fragileoverlays}) environment 
  \item enter the number of overlays as the first parameter to the 
  \texttt{overlays} environment  
  \item enter text contents with the \texttt{visible} command with the range
    of overlays showing this text content
\end{enumerate}

A simple example:

\begin{verbatim}
   % to be added in preamble
   \usepackage{overlays}
   ...
   \begin{overlays}{4}
   \begin{psli}[Title of slide]
   \begin{itemize}
     \visible{1-4}{\item first item of list} 
     \visible{2-4}{\item second list item}
     \visible{3-4}{\item 3rd list item}
     \visible{4}{\item final list item}
   \end{itemize}
   \end{psli}
   \end{overlays}
   ... 
\end{verbatim}

The following example uses the \texttt{alert} command to highlight lines
sequentially:

\begin{verbatim}
   \begin{overlays}{4}
   \begin{psli}[Title of slide]
   \begin{itemize}
     \alert{1}{\item first item of list} 
     \alert{2}{\item second list item}
     \alert{3}{\item 3rd list item}
     \alert{4}{\item final list item}
   \end{itemize}
   \end{psli}
   \end{overlays}
\end{verbatim}

The last example shows short text fragments which are shown sequentially using
the \texttt{only} command:

\begin{verbatim}
   \begin{overlays}{4}
     \begin{rsli}
       \only{1}{a short text, which will be replaced \ldots } 
       \only{2}{\dots by a second \ldots}
       \only{3}{\ldots and a third \ldots}
       \only{4}{\ldots and a final text.}
     \end{rsli}
   \end{overlays}
\end{verbatim}


The ``hidden'' text contents are written by \texttt{overlays} in the same
color as the background, default is white. If you use a different background
color, you have to change the color of the hidden text as well by assigning the
background color to the color name \texttt{background} (understood by the
\texttt{overlays} package).  In the following example you define a light
yellow as background:

\begin{verbatim}
   % (in the preamble)
   \definecolor{myyellow}{rgb}{0.96,0.98,0.72} % define color
   \definecolor{background}{named}{myyellow}   % color assigned to 
                                               % hidden text
   \pagecolor{myyellow}                        % color of slide background
\end{verbatim}



If you use one of the \emph{presentation schemes} described in section
\ref{sec:pres-scheme}, the necessary adjustments
for the background color will be made automatically.
For more details on \texttt{overlays}, see the documentation of the package.
\index{overlays|)}

%% hypertarget

Sometimes it is desirable to prepare \textbf{two or more consecutive
slides with the same page
number} independent
of the overlays package  (``manual overlays slides''). Therefore,
the page number of the current slide
can be set to the page number of the previous slide/page
with the command \verb+\pagenrconst+.
\index{pagenrconst@\verb+\pagenrconst+}

% \pagenrconst

\subsection{Add navigation symbols to slides}
\label{sec:NavigationSymbols}

\index{navigation symbols panel|(}
With the \texttt{navisymb} option of \texttt{elpres} a panel with navigation
symbols appears
in the right lower corner of pages:


{ \setlength{\parindent}{2em}
  \indent
    \texttt{$\ll$ $<$ $>$ $\gg$ $\leftarrow$ $\rightarrow$ [n]}
}

These commands ($\ll$: jump to the first page,
$<$: go to the previous page, $>$: go to the next page,
$\gg$: jump to the last page, $\leftarrow$:
go back in history, $\rightarrow$: go forward in history,
\texttt{[n]}: prompt for a
page number)
work with Adobe Acrobat Reader (Windows), and 
with \texttt{Okular} (Linux).
Some of these functions also work in presentation-mode
of the ``internal'' PDF-viewer in recent versions of the
\texttt{Firefox}-browser (Linux, Windows).
\index{navigation symbols panel|)}



\subsection{Add multimedia content to a presentation}
\label{sec:MultimediaContent}


% https://tex.stackexchange.com/questions/345431/how-to-include-multimedia-files-in-beamer

\index{multimedia files|(}
\index{audio files|see{multimedia files}}
\index{video files|see{multimedia files}}
This section describes inclusion of (links to) video and audio files in a 
presentation. With the
\verb+\auvimm{}{}+
\index{auvimm@\verb+\auvimm+}
command described below it is possible 
to lauch an external application for playing the video or audio
file. This will work in many situations for video and audio files
on Windows (Adobe Acrobat) and Linux (e.\,g. using
\index{Acrobat pdf viewer}%
\texttt{evince}%
\index{evince pdf file viewer@\texttt{evince} pdf file viewer})
systems in fullscreen mode.
Under certain conditions it is possible to define a ``poster area'' on a
slide and the pdf viewer starts (after clicking with the mouse pointer into
this area) the
presentation of the video within this frame.
The advantage of this technique is that it avoids starting
a window with the external application.
On Linux systems the \texttt{pdfpc}
\index{pdfpc viewer@\texttt{pdfpc} viewer}
viewer allows this, a detailed description
is shown below. In both situations, however, the multimedia file is not
incorprated into the pdf-file and must be available on the computer
at the time of the presentation, ideally in the same directory as the
\texttt{.pdf}-file.

The command \verb+\auvimm{text}{media-file}+ inserts \texttt{text} with a link
(``poster''), which allows to start the external default application for an
audio or video media file (\texttt{media-file}).  This worked flawlessly on
the computers of the author of this manual with Adobe Acrobat Reader (Windows)
and with \texttt{evince} (Linux).\footnote{The
  \texttt{\textbackslash auvimm} command uses the \texttt{\textbackslash href}
  command of the \texttt{hyperref} package}
An example for using this command:

\begin{verbatim}
   \auvimm{[sound-file]}{./audiofile.mp3}
\end{verbatim}

This produces the (text-based) ``button'' or ``poster'' \texttt{[sound-file]}.
You may also insert an image file as a poster with a graphics file using
\verb+\includegraphics+:

\begin{verbatim}
   \auvimm{\includegraphics{/path/to/poster.png}}{./audiofile.mp3}
\end{verbatim}

If you hit the poster area with the mouse pointer, it will launch the
default application for audio files in a separate window. This should
be possible even if the pdf-viewer works if fullscreen-mode at the time
of presentation.  If this does not work, it may be necessary to insert as
``launch'' command:

\begin{verbatim}
   \auvimm{[video-file]}{run:./videofile.mp4}
\end{verbatim}

This approach requires that you keep the presentation file
together with the multimedia files in the same folder, \textbf{also at
  the time of
  the presentation}. On an unknown computer system (where you have to show
your presentation) you should test this aspect of your presentation. Sometimes
a call to an external program might be blocked due security reasons in Acrobat
Reader.

%% begin comment out for tth  
   \fbox{\parbox{0.985\textwidth}{
%% end comment out for tth  
The \texttt{\textbackslash auvimm{}}-command in the current version
of \texttt{elpres} should be considered 
experimental.
Its behavior largely depends on the \texttt{pdf}-viewer and the software
used for playing audio/video files.
%% begin comment out for tth  
  }}
%% end comment out for tth  


On Linux systems, the \texttt{pdfpc}
\texttt{pdf}-viewer\footnote{\url{https://pdfpc.github.io/}: 
\texttt{pdfpc} --- A presenter console with multi-monitor support for PDF files} allows
presentation of a video embedded into a ``poster area''\footnote{e.\,g. defined
by a bitmap image with the same aspect ratio as the video}
on a slide
without starting the external viewer in a separate window. This is supported
with the \verb+\pdfpcmovie{}{}+-command from the
\texttt{pdfpc-movie}-package and with \verb+\auvimm{}{}+
using the ``\texttt{run:./video-file}'' launch-command. An example
with \verb+\auvimm{}{}+:

\begin{verbatim}
   \auvimm{\includegraphics[width=0.4\textwidth]{poster.png}}{run:./video.mp4}
\end{verbatim}

This \texttt{pdfpc}-specific extension seems to work only for video files,
but not for audio files. A possible solution is to convert an audio file into
a ``video format'', e.\,g. with the \texttt{ffmpeg} program.
\index{multimedia files|)}


\subsection{Prepare a ``handout'' from a presentation}

\label{secHandout}

In advance of a lecture it is often expected that you provide a ``handout''
of your presentation with more than one pages on a printed page. Therefore
you have several options, three of
them are described here.
If a presentation contains overlay-slides generated with the \texttt{overlays}
package as described in section \ref{secOverlay}, it should be recompiled
with the commands described in section \ref{sec:inact-overl-comm}
added to the preamble
in order to inactivate the overlays-specific commands.


\subsubsection{\LaTeX\ article document handout with a series of single-page pdf-files included}

One solution to this problem is (1) to generate a series of single-page
pdf-files and (2) to create a pdf\LaTeX\ document with the single
pages included.

\textbf{(1)} can be done with \texttt{gs} (the \texttt{ghostscript} program):
\index{ghostscript@\texttt{ghostscript}}

\begin{verbatim}
   gs -dNOPAUSE -dBATCH -sDEVICE=pdfwrite  -sOutputFile=page_%03d.pdf in.pdf
\end{verbatim}

with \texttt{in.pdf} as the initial presentation.  The syntax of \texttt{gs}
is described in detail on the ghostscript
website\footnote{\url{https://ghostscript.readthedocs.io/en/latest/Use.html}}.
On Windows systems the name of the ghostscript command may be
\texttt{gswin32c.exe} or  \texttt{gswin64c.exe}.
An alternative to \texttt{gs} is the \texttt{pdftk} command line tool:%
\index{pdftk@\texttt{pdftk}}%
\footnote{\url{https://www.pdflabs.com/tools/pdftk-the-pdf-toolkit/}}

\begin{verbatim}
    pdftk in.pdf  burst output  page_%03d.pdf
\end{verbatim}

The syntax of \texttt{pdftk} is explained in the documentation
(\texttt{pdftk~--help}).

\textbf{(2)} an example for a \LaTeX\ document which can serve as handout:

\begin{verbatim}
   \documentclass[12pt]{article}
   \usepackage[pdftex]{graphicx}
   \usepackage[latin1]{inputenc}
   \usepackage[a4paper,hmargin=2.4cm,top=24mm,bottom=28mm]{geometry}
   \newcommand{\PictScaleFact}{0.45}

   \begin{document}
   \centering
   \fbox{\includegraphics[width=\PictScaleFact\textwidth]{page_001.pdf}}
   \hspace{5mm}
   \fbox{\includegraphics[width=\PictScaleFact\textwidth]{page_002.pdf}}
   \\[2ex]
   \fbox{\includegraphics[width=\PictScaleFact\textwidth]{page_003.pdf}}
   \hspace{5mm}
   % ...
   \fbox{\includegraphics[width=\PictScaleFact\textwidth]{page_008.pdf}}
   % ...
   \end{document}
\end{verbatim}


\subsubsection{Create handout with the \texttt{pdfpages} package}

\label{sec:handout-pdfpages}

\index{handout|(}
Another, more comfortable option is to use the \texttt{pdfpages} package.
Here, it is not necessary to split up the presentation into single pages.
An example (in ``portait format'') which generates pages with 2x4 slides
(all slides except slide 2) using the
\verb+\includepdf[]{}+ command:

\begin{verbatim}
   \documentclass[12pt]{article}
   \usepackage[a4paper,hmargin=2cm,bottom=3.2cm]{geometry}
   \usepackage{pdfpages}
   \begin{document}
   
\includepdf[pages={1,3-last},nup=2x4,frame=true,%
   scale=0.78,%
   pagecommand={\thispagestyle{plain}}]{presentation.pdf}
   \end{document}
\end{verbatim}

Details on the optional parameters of \verb+\includepdf+ can
be found in the documentation of the \texttt{pdfpages} package. This code:

\begin{verbatim}
   \documentclass[12pt,landscape]{article}
   \usepackage[a4paper,hmargin=2cm,bottom=3.2cm]{geometry}
   \usepackage{pdfpages}
   \begin{document}
   
\includepdf[pages={1,3-last},nup=3x2,frame=true,%
   scale=0.78,%
   pagecommand={\thispagestyle{plain}}]{presentation.pdf}
   \end{document}
\end{verbatim}

generates pages in ``landscape format'' with 3x2 slides.

\subsubsection{Convert presentation pdf-file directly into handout file with \texttt{pdfjam}}

With the following command:

\begin{verbatim}
   pdfjam --nup 2x4 --frame true --scale 0.9 -o new.pdf in.pdf '1-4,6-22'   
\end{verbatim}

\begin{sloppypar}
\texttt{pdfjam}\footnote{which regerettably is only available on
  Linux or other Unix-like systems,
  on Windows systems, the procedure described
  above in section \ref{sec:handout-pdfpages} can be used as alternative}
creates a ``handout'' PDF document (\texttt{new.pdf}) fom
\texttt{in.pdf} with the slides nr. 1-4,6-22
arranged in  two columns and four rows.
With the additional option ``\texttt{--frame~true}'',
\texttt{pdfjam} draws a box
around each slide. More details can be found in the \texttt{pdfjam}
man page and the project
website\footnote{\url{https://github.com/DavidFirth/pdfjam/blob/master/README.md}}. 
\end{sloppypar}
\index{handout|)}



\subsection{Create presentations with hypertext elements}

\label{secHypertext}

\index{hypertext elements|(}
You may use two commands \verb+\fromlinktext{}{}+ and 
\verb+\totargettext{}{}+ for ``hypertext features''.  
As an example: as you normally will not 
insert \verb+\section{}+-like commands and therefore do not generate a ``
table of contents''-like page or a menu in \texttt{elpres}-presentations, 
you are able define links with:

\begin{verbatim}
   \totargettext{text}{link-label}
\end{verbatim}
% \hypertarget{target-name}{text}


which can be addressed by

\begin{verbatim}
   \fromlinktext{text}{link-label}
\end{verbatim}

% \hyperlink{target-name}{text}

If you use these commands to generate a \textbf{menu page} (as a hyperlinked
table of contents) with items pointing to
specific slides: 
\texttt{link-label} acts as label, \texttt{text} in the
\verb+\fromlinktext+ command is converted to a link (e.\,g. on a menu-page).
If you click on this text ``button'', you jump to
the \texttt{text} labelled with
\verb+\totargettext+; \texttt{link-label} has to be identical in a
\verb+\fromlinktext+ and \verb+\totargettext+ pair,
\texttt{text} in the link
and the target, of course need not to be identical.

The default for the link border color is red, it may be changed with
\verb+\hypersetup{}+, a command from the \texttt{hyperref}
package\footnote{both, \texttt{\textbackslash totargettext} and
\texttt{\textbackslash fromlinktext} are
are redefinitions of commands from the \texttt{hyperref} package}.
The option for the color of the frame around the link is
\texttt{linkbordercolor},
colors must be defined whith three figures \texttt{[0..1]} according to the RGB
color model:

\begin{verbatim}
   \hypersetup{linkbordercolor={0.6 0.6 0.6}}
\end{verbatim}

This defines a light gray color for the link border, more details on
\verb+\hypersetup{}+ are
described in the documentation of the \texttt{hyperref} package.
\index{hypertext elements|)}


\subsection{Fill background of a presentation with bitmaps}

\label{secBackgroundWallpaper}

\subsubsection{\texttt{Wallpaper} package}

\index{wallpaper background|(}
To create a slide background with a graphical wallpaper 
background using bitmap files you may use the \texttt{wallpaper}
package\footnote{written by Michael H.F. Wilkinson and available on CTAN}. 
Load the \texttt{wallpaper} package with

\begin{verbatim}
   \usepackage{wallpaper}
\end{verbatim}

\begin{sloppypar}
in the preamble. In order to generate a background based on bitmap
file \texttt{background.png}, enter 
\end{sloppypar}

\begin{verbatim}
   \CenterWallPaper{1}{background.png}
\end{verbatim}

before the contents of the presentation\footnote{i.\,e. following
\textbackslash begin\{document\}}. 
This works best with bitmaps with an appropriate aspect ratio, in the case of
an 4x3 screen format a bitmap picture of 640x480 pixel would fit perfectly. 
Moreover bitmap files may be
used as tiles as described in the \texttt{wallpaper} documentation like

\begin{verbatim}
   \TileSquareWallPaper{4}{background.png}
\end{verbatim}

More details on this topic may be found in the \texttt{wallpaper}
documentation.

\subsubsection{\texttt{Eso-pic} package}
\label{secBackgroundEsopic}

Another package which allows you to paint the background with a picture is
\texttt{eso-pic}\footnote{written by Rolf Niepraschk and available on CTAN}:

\begin{verbatim}
   \usepackage{eso-pic}
   
   ...

   \AddToShipoutPictureBG*{
      \includegraphics[height=\paperheight]{background.png}
   }
\end{verbatim}

\verb+\AddToShipoutPictureBG{}+ puts the picture on every page, 
\verb+\AddToShipoutPictureBG*{}+ puts it on to the current page,
\verb+\ClearShipoutPictureBG+ clears the background beginning with the current
page.
Details of \texttt{eso-pic}'s commands can be found in the documentation.
\index{wallpaper background|)}


\subsection{Put color to your presentation: presentation schemes}
\label{sec:pres-scheme}

\index{presentation schemes|(}
\index{color schemes|see{presentation schemes}}
\index{schemes, color|see{presentation schemes}}

\texttt{Elpres} provides predefined \textbf{presentation schemes}: color and
layout schemes which are applied to a presentation with
\verb+\usepackage{}+. Their naming follows the convention
\mbox{\texttt{elpres...scheme}}, an example is
\texttt{elpreswhiteredscheme}.\footnote{the name of the corresponding style
or package file is \texttt{elpreswhiteredscheme.sty}} These
\texttt{elpres}-schemes also support incremental slides with the
\texttt{overlays} package. With an activated presentation scheme, the elpres
class option \texttt{ballsymb} changes default itemize symbols to
appropriately colored ``ball'' symbols. Moreover, these schemes define color
names which may be used by the author of a presentation:\footnote{Color
names for use by the presentation author start with \texttt{ep...}}

\begin{description}%[noitemsep]
\item[\texttt{eptextcolor}] ``normal text color'' (default values for
  documents without
  presentation scheme included in \texttt{elpres.cls}: black)
\item[\texttt{ephighlightcolor}] color, which may be used for highlighting
  text (default: dark red)
\item[\texttt{eptitlecolor}] color of the title of slides, may be used to
  highlight the title in the title slide (default: dark red)
\end{description}

These presentation schemes available in the current version of \texttt{elpres}:

\begin{sloppypar}

\begin{description}%[noitemsep]

% \item[\texttt{elpresplainscheme}] ..

\item[\texttt{elpresbluelightgrayscheme}] dark blue background of slides;
  \texttt{eptextcolor}: light gray;
  \texttt{eptitlecolor}: turquoise;
  \texttt{ephighlightcolor}: yellow;
  symbols in \emph{itemize} environment: pale blue bullets;
  highlighted elements in \emph{enumerate} and \emph{description}
  environments: pale blue
  
\item [\texttt{elpresgrayscheme}] gray background of slides;
  \texttt{eptextcolor}: dark gray;
  \texttt{eptitlecolor}:
  dark green; \texttt{ephighlightcolor}: dark red;
  symbols in \emph{itemize} environment: gray bullets; highlighted elements
  in \emph{enumerate} and \emph{description} environments: dark green

\item [\texttt{elpreswhitebluescheme}] white background of slides;
  \texttt{eptextcolor}: blue-black;
  \texttt{eptitlecolor}: blue;
  \texttt{ephighlightcolor}: dark red;
  symbols in \emph{itemize} environment: blue bullets; highlighted elements
  in \emph{enumerate} and \emph{description} environments: blue

\item [\texttt{elpreswhiteredscheme}] white background of slides;
  \texttt{eptextcolor}: red/brown-black;
  \texttt{eptitlecolor}: dark red;
  \texttt{ephighlightcolor}: dark blue;
  symbols in \emph{itemize} environment: red bullets;
  highlighted elements in \emph{enumerate} and \emph{description}
  environments: red

\item [\texttt{elpreswhitetealscheme}] white background of slides;
  \texttt{eptextcolor}: black-teal;
  \texttt{eptitlecolor}: teal;
  \texttt{ephighlightcolor}: dark red;
  symbols in \emph{itemize} environment: blue bullets; highlighted elements
  in \emph{enumerate} and \emph{description} environments: dark teal
  
\end{description}
\end{sloppypar}

In all of these schemes, the colors of borders for internal and external
links have been adjusted.
\index{presentation schemes|)}


% Die Farben in LibreOffice / rgb (266) / rgb (1)
% ----------------------------------------------
% Tango himmelblau 1 / 114 159 207 /
% Rot 5 / 153 0 102 /
% Grau 6 / 128 128 128 /

\section{Recipes for use of \texttt{elpres}}
\label{sec:elpres-recipes}

\index{recipes for elpres presentations@recipes for \texttt{elpres} presentations}

\subsection{Colored text boxes}
\label{sec:colored-text-boxes}

\index{text boxes, colored}
Colored text boxes often help the presentation author to highlight small text
fragments and to make slides more ``readable'' and more attractive.

\subsubsection{Text boxes with \texttt{colorbox}, \texttt{fcolorbox}}
\label{sec:simple-text-boxes}


A simple method for colored text boxes uses the \texttt{minipage} environment
and the \texttt{colorbox} or \texttt{fcolorbox}
commands (\texttt{xcolor} package):

\begin{verbatim}
   \begin{psli}[Slide with a textbox]
   \begin{center}
     \colorbox{eptitlecolor}{\begin{minipage}{0.85\textwidth}
        \centering \color{white} \vspace{2ex}
        {\Large
        Text centered in a colorbox
        } \par
       \vspace{2ex}
     \end{minipage}}
   \end{center}
   \end{psli}
\end{verbatim}

The advantage of these simple colored text boxes: they  are compatible
with incremental slides using the \texttt{overlays} package.

\subsubsection{Text boxes with the \texttt{tcolorbox} package}
\label{sec:tcolorbox-boxes}

The \texttt{tcolorbox} package provides great support for colored boxes. A
minimal example:

\begin{verbatim}
   \begin{psli}[Slide with a textbox]
    \begin{tcolorbox}[colframe=eptitlecolor!95!black]
     \textcolor{eptitlecolor}{This is a textbox generated with the
       \texttt{tcolorbox} package}
     \tcblower This the lower part of the tcolorbox
    \end{tcolorbox}  
   \end{psli}
\end{verbatim}

An example with a titled colored textbox:

\begin{verbatim}
   \begin{psli}[Slide with a textbox]
    \begin{tcolorbox}[colback=eptitlecolor!10!white,%
       colframe=eptitlecolor!95!black,%
       title=Heading of a textbox]
       This is another \textbf{tcolorbox}.
     \tcblower
        Here, the lower part of the box.
      \end{tcolorbox}
   \end{psli}
   
\end{verbatim}

The manual of the \texttt{tcolorbox} provides detailed instructions for the 
use of this package.  
However, incremental slides using the \texttt{overlays} package may not
work with color boxes generated by the \texttt{tcolorbox} package.  These
problems are avoided with the colored text boxes  described in section
\ref{sec:simple-text-boxes}.

\subsection{Alternative format for footnotes on slides}
\label{sec:AltFootnoteFormat}


Footnotes on slides may be inserted with \LaTeX's \verb+\footnote{}+
command.
However, the layout of footnotes generated by \verb+\footnote{}+ 
is not adequate
for small comments or additions to the main text on slides.
The following definition of an alternative footnote command 
(\verb+\altfnarea{}+)%  
\index{alternative footnote command}
\index{footnote command, alternative}
\index{altnfarea@\verb+\altfnarea+}
requires that the 
\texttt{eso-pic} package is loaded:

\begin{verbatim}
   \usepackage{eso-pic}

   \newcommand{\altfnarea}[1]{%
   \AddToShipoutPictureFG*{
   \setlength{\unitlength}{1mm}
       \put(1.2,1.2){\parbox[b]{0.9\paperwidth}%
       {\scriptsize \textcolor{darkgray}{#1}}}}}
   
\end{verbatim}

The definition of \verb+\altfnarea+ is currently not included in
\texttt{elpres.cls}, so you should copy this definition of the command
into the preample of your presentation document. Please take care to
include the \texttt{eso-pic} package before the \texttt{altfnarea}
definition. You may experiment with different values in this
\verb+\newcommand{}{}+ definition.  The command can be used within
the code of a slide the following way:

   \begin{verbatim}
   \begin{psli}[Title of a slide demonstrating alternative 
      footnotes\textsuperscript{1}]
   \begin{itemize}
   \item first itemized line\textsuperscript{2}
   \item second itemized line
   \end{itemize}
   \altfnarea{\textsuperscript{1}Addition to the title, 
      \textsuperscript{2}Explanation of the first line}
   \end{psli}
   \end{verbatim}

Very often it will be advisable \emph{not} to use footnotes in slides.
However, speakers participating in scientific conferences usually are
expected to include sources of figures, tables etc. on their slides.


\subsection{Inactivate \texttt{overlays} commands}
\label{sec:inact-overl-comm}

\begin{sloppypar}
If a pesentation has been prepared with overlay slides (using the
\texttt{overlays}) package, the \LaTeX{} source code will contain
\verb+\visible+, \verb+\only+ and \verb+\alert+ commands and
\texttt{overlays} and \texttt{fragileoverlays} environments.

Sometimes, it is desirable to generate a version of such a presentation
without overlays effects, for example for preparation of a printed handout
(see section \ref{secHandout}).
The effect of these overlay-specific instuctions can be
inactivated with these commands:
\end{sloppypar} 


\begin{verbatim}
   \overlaysoff
   \alertsoff
\end{verbatim}

which should be written into the preamble following
\verb+\usepackage{overlays}+.  This works with v2.12 (and probably later
versions) of \texttt{overlays}.  With earlier versions you may copy the
commands

\begin{verbatim}
   \renewenvironment{overlays}[1]{}{}%
   \renewenvironment{fragileoverlays}[1]{}{}%
   \renewcommand{\visible}[2]{#2}%
   \renewcommand{\only}[2]{#2}%
   \renewcommand{\alert}[2]{#2}%
\end{verbatim}

into the preamble following \verb+\usepackage{overlays}+, however,
inactivation of \verb+\only+ by this method
has a different effect, so it is recommended to
update the \texttt{overlays} package.

\subsection{Convert an \texttt{elpres}-presentation (PDF) into an Impress (LibreOffice) or Powerpoint presentation}
\label{sec:convert-into-libreoff}

\subsubsection{PDF \ldots\ PNG \ldots\ Impress/Powerpoint}

\index{Powerpoint presentation} \index{Impress (LibreOffice) presentation}
Sometimes it is unclear, if organizers of a (scientific) meeting allow
presentations with \texttt{.pdf}-files or if they insist on files in Powerpoint
format. This is an unpleasant situation for a lecturer, but in such a situation
it is better, to be prepared.

There is no elegant way for a reliable solution.  First, you may use one
of the free or commercial online conversion tools (intended to convert
\texttt{.pdf} into \texttt{.ppt} or \texttt{.pptx}-files). Results are
often not satisfactory.

An alternative, more reliable way preferred
by the author of this manual:


\begin{enumerate}
\item convert the \texttt{.pdf}-file of the presentation into a series of bitmap
  image files, (e.\,g. in \texttt{.png} format)
\item import these image \texttt{.png} files into ``empty'' slides
  of an Impress
  (LibreOffice or OpenOffice) presentation and save the presentation in
  \texttt{.odp}, \texttt{.ppt} or \texttt{.pptx}-Format
\end{enumerate}

\textbf{Details on step 1 --- generate \texttt{.png} images:}

Please copy the following two
commands (please ignore the line-breaks in the lines
beginning with ``\texttt{gs}'')\footnote{\i.\,e. ``\texttt{gs ... input.pdf}'' should
  be entered as one line and ``\texttt{gs ... temp.pdf}'' should be entered as a
  second line of a shell script.}
into a short shell
script (Windows: a \texttt{.bat} or \texttt{.cmd} script) and ``run'' or
execute the shell script:



\begin{verbatim}
   gs -dNOPAUSE  -dBATCH -sDEVICE=pdfwrite  -dDEVICEWIDTHPOINTS=793.7 
      -dDEVICEHEIGHTPOINTS=595.3 -dPDFFitPage -sOutputFile=temp.pdf input.pdf

   gs -dNOPAUSE -dBATCH -sDEVICE=png256 -r300 -dTextAlphaBits=4 
      -sOutputFile=page_%03d.png temp.pdf
\end{verbatim}

\begin{sloppypar}
  Please do not forget to adjust the name of the input file
  (\texttt{input.pdf} in the script above).
  The first command adjusts the page size of the \texttt{.pdf}-file to
  \texttt{28~x~21~cm} which is written into the temporary file
  \texttt{temp.pdf}:
  this format avoids the necessity to resize of the images in LibreOffice 
  (for other slide formats see the footnote below\footnote{%
     The values assigned to \texttt{-dDEVICEWIDTHPOINTS} and 
     \texttt{-dDEVICEHEIGHTPOINTS}
     are valid for a width:height aspect ratio of 4:3 assuming a
     presentation page size (LibreOffice) of \texttt{28cm~x~21cm}:
     \texttt{(28/2.54)*72=793.7pt},
     \texttt{(21/2.54)*72=595.3pt}. Values for width:height aspect ratios of
     16:9 \texttt{(28cm~x~15.75cm)} are \texttt{793.7/446.46} and for a ratio
     of 16:10 \texttt{(28cm~x~17.5cm)}:  \texttt{793.7/496.1}.
  }).
  The second command generates a series of \texttt{.png}-files:
  \texttt{page\_001.png}, \texttt{page\_002.png} ... \texttt{page\_0nn.png}.
  Instead of the \texttt{png256} driver, \texttt{png16m} can be used.
  On Windows systems the name of the ghostscript command is
  \texttt{gswin32c} or \texttt{gswin64c} (instead of \texttt{gs}).
\end{sloppypar}

If you wish a higher resolution of the \texttt{.png} images,
change \texttt{-r300} to
\texttt{-r600}.
% \begin{verbatim}
%    gs -dNOPAUSE -dBATCH -sDEVICE=png256 -r600 
%            -dTextAlphaBits=4 -sOutputFile=page_%03d.png input.pdf
% \end{verbatim}
The option  \mbox{\texttt{-dTextAlphaBits=4}} improves font
antialiasing.\footnote{Details can be found in the documentation of 
\texttt{gs}  \url{https://ghostscript.readthedocs.io/en/latest/Use.html}}

\textbf{Details on step 2 --- Import \texttt{.png} images:}

Open an empty presentation with LibreOffice Impress, adjust the appropriate
slide format (e.\,g. 4:3), import the \texttt{.png} files with
\emph{Insert~/ Image}
% , then adjust the size of the included image with
% \emph{Format~/ Object~and~Shape~/ Position~and~Size (F4)} in the
% \emph{Position~and~Size} dialog box please enter: \emph{Postion~X}:
% \texttt{0.0cm}, \emph{Postion~Y}: \texttt{0.0cm}, \emph{Width}:
% \texttt{28.0cm}, activate the checkbox \emph{Keep~ratio}
into consecutive empty slides.\footnote{Menu
  structure of LibreOffice v6.4, v7.3}. The presentation can now be saved
in the format of Impresss (\texttt{.odp}) or
one of the Microsoft Powerpoint formats (\texttt{.ppt} or \texttt{.pptx}).

\subsubsection{PDF \ldots\ SVG \ldots\ Impress/Powerpoint}

\label{sec:PdfSvgImpress}

An alternative way to make a Powerpoint/Impress presentation from the
.pdf-file is to convert it into a series of \textbf{.pdf}-pages, convert
these into \texttt{.svg} format and to import the \texttt{.svg}-files into
the Impress/Powerpoint presentation. The first step:%
\footnote{The command in a Windows terminal should be:
  \texttt{pdftk input.pdf burst output p\_\%\%02d.pdf}}

\begin{verbatim}
   pdftk input.pdf burst output p_%02d.pdf
\end{verbatim}

converts the file \texttt{input.pdf} into a series of files
\texttt{p\_01.pdf}, \texttt{p\_02.pdf}, \texttt{p\_03.pdf}, \ldots, which
can then be converted with the command \texttt{pdftocairo}:%
\footnote{Available in the \texttt{poppler-utils} package, a Windows version
may be obtained at: \url{https://github.com/oschwartz10612/poppler-windows}}


\begin{verbatim}
   pdftocairo -svg p_01.pdf p_01.svg 
   pdftocairo -svg p_02.pdf p_02.svg
   pdftocairo -svg p_03.pdf p_03.svg
   ...
\end{verbatim}

This series of commands is best executed from a shell script. The resulting 
\texttt{.svg} files are then imported into an Impress/Powerpoint presentation.
 
\subsection{The \texttt{pdfpc}-viewer: add notes}

\index{pdfpc viewer@\texttt{pdfpc} viewer}
\index{notes for slides, pdfpc viewer@notes for slides, \texttt{pdfpc} viewer}
The pdf-vierwer \texttt{pdfpc} for
Linux\footnote{\url{https://pdfpc.github.io/}}
(see section \ref{sec:MultimediaContent})
is a comfortable program which
allows output of presentations to two monitors or a monitor and a
beamer/projector.  The normal presentation window is assigned to the
beamer/projector, the other output for the presenter shows the actual slide,
the next slide and an area for notes. As you are viewing a certain slide,
press \texttt{n} to start editing a note. The editing mode is stopped with
\texttt{Esc}.  The text of these notes is stored in the file
\texttt{filename.pdfpc}.  Further information can be found in the \texttt{man}
page. 


\section{License}

This class is distributed under the \textsl{\LaTeX{} Project Public License}
(LPPL) which may be downloaded from
\url{http://www.latex-project.org/lppl.txt}. No warranty is provided for this
work (as stated in the LPPL).



\section{Version history}

\textbf{v0.1} (19.6.2004): initial version.  \textbf{v0.2} (1.9.2004): page
numbers now changed to footnotesize, left and right margins slightly changed,
\texttt{cenumerate} and \texttt{cdescription} environments added.
\textbf{v0.2a} (19.9.2004): Section ``License'' added to the documentation.
\textbf{v0.2b} (17.10.2004): Documentation completed: description of the
\verb+\distance{}+ command included.  \textbf{v0.2c} (28.11.2004):
Documentation completed (section \ref{secHypertext} added).  \textbf{v0.2d}
(25.12.2004): Documentation completed (section \ref{secBackgroundWallpaper}
added).  \textbf{v0.2e} (15.04.2005): Documentation completed (sections
\ref{secBackgroundEsopic} and \ref{secHandout} added).  \textbf{v0.3}
(12.08.2005): new (class) options for font selection: \texttt{tmrfont} (Times
Roman), \texttt{helvetfont} (Helvetica), \texttt{cmfont} (Computer Modern),
\texttt{sansfont} (Sans Serif: default).  Documentation updated, sections
\ref{secUsageGraph} and \ref{secTwoColumns} added.  \textbf{v0.4}
(20.01.2018): New class options for different screen aspect ratios
\texttt{4x3}, \texttt{16x9}, \texttt{16x10}; ``compressed'' list environments
modified; documentation completed: packages for use with \texttt{elpres}:
\texttt{enumitem} (alternative list environments), \texttt{overlays} (overlay
support: incremental slides); section \ref{secHandout} was completely
rewritten.  \textbf{v0.4a} (24.01.2018): Documentation completed.
\textbf{v0.5} (12.07.2020): New class options: \texttt{nofonts} (no font
selection) and \texttt{navisymb} (inserts a panel of symbols for navigation),
new commands: \verb+auvimm+ (generates a link to external audio, video files),
\verb+\fromlinktext+ and \verb+\totargettext+ (create links within a
presentation).  Documentation has been partly rewritten.  \textbf{v0.6}
(19.08.2020): Access to the \texttt{fancyhdr} ``header fields''
(\verb+\lhead{}+, \verb+\chead{}+ and \verb+\rhead{}+) is now explicitly
blocked.  New command \verb+\slidetitlecolor{}+ for the text color of titles
in \texttt{psli}-slides.  New: style files for presentation schemes added.
Documentation has been significantly completed: sections
\ref{sec:footerlayout}, \ref{sec:use-colors}, \ref{sec:pres-scheme} and
\ref{sec:elpres-recipes} added, section \ref{sec:installation} on installation
rewritten.  \textbf{v0.7} (20.02.2021): parts of the manual (this file) have
been rewritten and an index has been included. The command \verb+\pagenrconst+
has been added.  \textbf{v0.8} (28.02.2021): Installation instructions have
been updated in section \ref{sec:Introduction} and detailed instructions were
added to \texttt{README.md}. Section \ref{sec:inact-overl-comm} was updated. A
bug in the  \texttt{elpres...scheme.sty} files was fixed (incorrect value for
\verb+\labelitemiv+).
\textbf{v0.9} (01.08.2021): New class option \texttt{bulletsymb} added, an
archive \texttt{elpres.tds.zip} is added for \texttt{ctan.org} upload.
\textbf{v1.0}: Class option \texttt{ballsymb} added, manual updated.
\textbf{v1.0.1}: Corrected upload to \texttt{ctan.org} 
\textbf{v1.1} (28.05.2024): Fixes and corrections in \texttt{elpres.cls};
manual: parts of section \ref{secUsage} have been rewritten,
section \ref{sec:convert-into-libreoff} has been completed: new section
\ref{sec:PdfSvgImpress}.
Definition of a an alternative command for footnotes (\verb+\altfnarea+) 
has been added to the manual (section \ref{sec:AltFootnoteFormat}).


\printindex

\end{document}

